

\section{Hardness of LWE with Binary Secret}
\label{sec:shortsecret}



The following is the main theorem of this section. 

%\dnote{I changed~$\alpha, \delta \in (0,1)$ to~$\alpha,\delta >0$}

\begin{theorem}\label{thm:zolwe}
Let~$k, q \ge 1$, and $m \ge n \ge 1$ be integers, and let $\epsilon \in (0,1/2)$, $\alpha, \delta >0$,
be such that $n \ge (k+1) \log_2 q + 2 \log_2 (1/\delta)$,
$\alpha \ge \sqrt{\ln(2n(1+1/\eps))/\pi}/q$.
There exist three (transformation) reductions from $\lwe_{k, m, q, \alpha}$ to $\zolwe_{n, m, q, \le \sqrt{10n} \alpha}$,
such that for any algorithm for the latter problem with advantage $\zeta$, at least one of the reductions produces 
an algorithm for the former problem with advantage at least 
\begin{equation}\label{eq:advantagelossinzolwe}
(\zeta-\delta)/(3m) - 41 \eps/2 - \sum_{p|q,~p~\text{prime}} p^{-k-1}\,.
\end{equation}
\end{theorem}

By combining Theorem~\ref{thm:zolwe} with the reduction in Corollary~\ref{cor:mod-reduction} (and noting that $\{0,1\}^n$ is $(\sqrt{n},0)$ bounded), we can replace the $\zolwe$ problem above with $\zolwe_{n,m,q',\beta}$ for any $q' \ge 1$ and $\xi>0$ where
\[
\beta := \left(10 n \alpha^2 + \frac{4n}{\pi q'^2} \ln(2n(1+1/\xi))\right)^{1/2},
\]
while decreasing the advantage in~\eqref{eq:advantagelossinzolwe} by $14 \xi m$. Recalling that 
$\lwe$ of dimension $k=\sqrt{n}$ and modulus $q=2^{k/2}$ (assume $k$ is even) is known to be 
classically as hard as $\sqrt{n}$-dimensional
lattice problems (Theorems~\ref{thm:worstcaseaveragecase} and~\ref{thm:searchtodecisionmicpei}), 
this gives a formal statement of Theorem~\ref{thm:informal}. The modulus $q'$ can be taken
almost as small as $\sqrt{n}$.

For most purposes the sum over prime factors of $q$ in~\eqref{eq:advantagelossinzolwe} is
negligible. For instance, in deriving the formal statement of Theorem~\ref{thm:informal}
above, we used a $q$ that is a power of $2$, in which case the sum is $2^{-k-1}=2^{-\sqrt{n}-1}$, which
is negligible. If needed, one can improve this by applying the modulus switching reduction (Corollary~\ref{cor:mod-reduction-2})
before applying Theorem~\ref{thm:zolwe} in order to make $q$ prime. (Strictly speaking, one also needs to apply Lemma~\ref{lem:lelwe} to replace the ``unknown noise'' variant of $\lwe$ given by Corollary~\ref{cor:mod-reduction-2} with the fixed noise variant.) This improves the advantage loss
to $q^{-\sqrt{n}-1}$ which is roughly $2^{-n}$. 

In a high level, the proof of the theorem follows by combining three main steps. The first, given in
Section~\ref{sec:firstiserrorless}, reduces $\lwe$ to a variant in which the first
equation is errorless. The second, given in Section~\ref{sec:extlwe}, reduces
the latter to the intermediate problem $\extlwe$, another variant of $\lwe$ in which some information 
on the noise elements is leaked. Finally, in Section~\ref{sec:theshortsecretreduction},
we reduce $\extlwe$ to $\lwe$ with $\{0,1\}$ secret. We note that the first reduction is
relatively standard; it is the other two that we consider as the main contribution
of this section. We now proceed with more details (see also Figure~\ref{fig:summary}). 

\begin{proof}
First, since $m \ge n$, Lemma~\ref{lem:firsterrorlesshard} provides a transformation reduction
from $\lwe_{k, m, q, \alpha}$ to first-is-errorless $\lwe_{k+1, n, q, \alpha}$, while reducing the advantage by at most $2^{-k+1}$. 
Next, Lemma~\ref{lem:redtoextlwe} with $\cZ = \{0,1\}^n$, which is of quality $\xi=2$ by 
Claim~\ref{clm:qualityofzeroone}, reduces the latter problem to $\extlwe_{k+1, n, q,\sqrt{5}\alpha,\{0,1\}^n}$ while 
reducing the advantage by at most~$33 \eps/2$. Then, Lemma~\ref{lem:multiinstanceextlwe} reduces the latter problem
to $\extlwe^m_{k+1, n, q, \sqrt{5} \alpha,\{0,1\}^n}$, while losing a factor of~$m$ in the advantage. 
Finally, Lemma~\ref{lem:zolwe} provides three reductions to $\zolwe_{n, m, q, \le \sqrt{10n} \alpha}$: two from the latter problem, and one from
$\lwe_{k+1, m, q,\sqrt{5n} \alpha}$, guaranteeing that the sum of advantages is at least the original advantage minus $4 m \eps + \delta$.
Together with the trivial reduction from $\lwe_{k, m, q,\alpha}$ to $\lwe_{k+1, m, q, \sqrt{5n} \alpha}$ (which incurs no loss in advantage), this
completes the proof. 
\end{proof}

\begin{figure}[ht]%
\centering
\begin{tikzpicture}[node distance=1cm, auto]  
\tikzset{
    mynode/.style={rectangle,draw=black, top color=white, bottom color=yellow!5,thick, inner sep=0.5em, minimum height=2em, text centered},
    myarrow/.style={->, >=latex', shorten >=1pt, thick},
    mylabel/.style={text width=7em, text centered} 
}  
%\node[mynode] (final) {$\zolwe_{n, m',q, \sqrt{10n} \alpha}$};  
%\node[mynode, above=0.5cm of final] (atmost) {$\zolwe_{n, m,q, \le \sqrt{10n} \alpha}$}; 
\node[mynode] (atmost) {$\zolwe_{n, m,q, \le \sqrt{10n} \alpha}$}; 
\node[mynode, above right=1.2cm of atmost] (second) {$\lwe_{k+1, m, q, \sqrt{5n} \alpha}$};
\node[mynode, above left=1.2cm of atmost] (first) {$\extlwe^m_{k+1, n, q,\sqrt{5} \alpha,\{0,1\}^n}$};  
%\node[mynode, right=of second] (third) {discrete $\lwe^n_{k, m, q,\beta}$};
%\node[mynode, above=0.5cm of third] (thirdabove) {discrete $\lwe_{k, m, q,\beta}$};
%\node[mynode, above=0.5cm of thirdabove] (thirdaboveabove) {$\lwe_{k, m, q,\beta}$};
\node[mynode, above=0.5cm of first] (firstabove) {$\extlwe_{k+1, n, q, \sqrt{5} \alpha,\{0,1\}^n}$};
\node[mynode, above=0.5cm of firstabove] (firstaboveabove) {1$^{\mbox{\scriptsize{st}}}$ errorless $\lwe_{k+1, n,q,\alpha}$};
\node[mynode, above=5cm of atmost] (firstaboveaboveabove) {$\lwe_{k, m, q,\alpha}$};

%\draw[myarrow] (atmost)  --  (final);	
\draw[myarrow] (first)  --  (atmost);	
\draw[myarrow] (second)  --  (atmost);	
%\draw[myarrow] (third)  --  (atmost);	
%\draw[myarrow] (thirdabove)  --  (third);	
%\draw[myarrow] (thirdaboveabove)  --  (thirdabove);	
\draw[myarrow] (firstabove)  --  (first);	
\draw[myarrow] (firstaboveabove)  --  (firstabove);	
\draw[myarrow] (firstaboveaboveabove)  --  (firstaboveabove);	
\draw[myarrow] (firstaboveaboveabove)  --  (second);	
%\draw[myarrow] (firstaboveaboveabove)  --  (thirdaboveabove);	

%\node[mylabel, above left=0cm of final] (label1) {Lemma~\ref{lem:lelwe}};  
%\node[mylabel, below left=0cm of first] (label2) {Lemma~\ref{lem:zolwe}};  
\node[mylabel, above left=0cm of first] (label3) {Lemma~\ref{lem:multiinstanceextlwe}};  
\node[mylabel, above left=0cm of firstabove] (label4) {Lemma~\ref{lem:redtoextlwe}};  
\node[mylabel, below left=0cm of firstaboveaboveabove.west] (label5) {Lemma~\ref{lem:firsterrorlesshard}};  
%\node[mylabel, above right=0cm of third] (label6) {Lemma~\ref{lem:multiinstancelwe}};  
%\node[mylabel, above right=0cm of thirdabove] (label7) {Lemma~\ref{lem:conttodiscrete}};  
\node[rectangle,draw=black, thick, inner sep=0.5em, minimum width=5cm, fill=blue!10!white, fill opacity=0.8, text opacity=1, dashed, above=0.3cm of atmost] (emphasize) {Lemma~\ref{lem:zolwe}};
\end{tikzpicture} 
\caption{Summary of reductions used in Theorem~\ref{thm:zolwe}}%
\label{fig:summary}%
\end{figure}




\subsection{First-is-errorless LWE}
\label{sec:firstiserrorless}

We first define a variant of LWE in which the first equation is given without error, and then show in
Lemma~\ref{lem:firsterrorlesshard} that it is still hard. 

\begin{definition}\label{def:firsterrorless}
For integers $n,q \ge 1$ and an error distribution $\phi$ over $\R$, the ``first-is-errorless'' variant of the $\lwe$ problem
is to distinguish between the following two scenarios.
In the first, the first sample is uniform over $\T_q^n \times \T_q^{}$ and the rest are uniform
over $\T_q^n \times \T$. In the second, there is an unknown uniformly distributed $\vecs \in \{0,\ldots,q-1\}^n$,
the first sample we get is from $A_{q,\vec{s},\{0\}}$ (where $\{0\}$ denotes the distribution that is deterministically zero)
and the rest are from $A_{q,\vec{s},\phi}$.
\end{definition}


\begin{lemma}\label{lem:firsterrorlesshard}
For any $n \ge 2$, $m, q \ge 1$, and error distribution $\phi$,
there is an efficient (transformation) reduction from $\lwe_{n-1,m,q,\phi}$ to the first-is-errorless variant of $\lwe_{n,m,q,\phi}$
that reduces the advantage by at most $\sum_p p^{-n}$, with the sum going over all prime factors of $q$.
\end{lemma}
Notice that if $q$ is prime the loss in advantage is at most $q^{-n}$. Alternatively, for any number $q$ we can bound it by
\[
\sum_{k \ge 2} k^{-n} \le 2^{-n} + \int_{2}^\infty t^{-n} \mathrm{d} t \le 2^{-n+2},
\]
which might be good enough when $n$ is large.
\begin{proof}
The reduction starts by choosing a
  vector $\veca'$ uniformly at random from
  $\{0,\ldots,q-1\}^n$. Let $r$ be the greatest common divisor of the coordinates of $\veca'$. 
	If it is not coprime to $q$, we abort. The probability that this
  happens is at most 
\[
 \sum_{p~\text{prime}, \ p|q} p^{-n}.
\]
Assuming we do not abort, we proceed by finding a matrix $\matU \in \Z^{n \times n}$ that is invertible modulo $q$ and
whose leftmost column is $\veca'$.
Such a matrix exists, and can be found efficiently. For instance, using the extended GCD algorithm, we find an $n\times n$ unimodular matrix
$\matR$ such that $\matR \veca' = (r,0,\ldots,0)^T$. Then $\matR^{-1} \cdot \mathrm{diag}(r,1,\ldots,1)$ is the desired matrix.
We also pick a uniform element $s_0 \in \{0,\ldots,q-1\}$. The reduction now proceeds as follows. The first sample it outputs is
$(\veca'/q, s_0/q)$. The remaining samples are produced by taking a sample $(\veca,b)$ from the given oracle, picking
a fresh uniformly random $d \in \T_q$, and outputting $(\matU(d|\veca),b+(s_0 \cdot d))$ with the vertical bar denoting concatenation.
It is easy to verify correctness: given uniform samples, the reduction outputs uniform samples
(with the first sample's $b$ component uniform over $\T_q$), up to statistical distance $2^{-n+1}$;
and given samples from $A_{q,\vec{s},\phi}$, the reduction outputs one sample from $A_{q,\vecs',\{0\}}$
and the remaining samples from $A_{q,\vecs',\phi}$, up to statistical distance $2^{-n+1}$,
where $\vecs' = (\matU^{-1})^T (s_0|\vecs) \bmod q$.
This proves correctness since $\matU$, being invertible modulo $q$, induces a bijection on $\Z_q^n$, and so $\vecs'$ is uniform in $\{0,\ldots,q-1\}^n$.
\end{proof}

\subsection{Extended LWE}
\label{sec:extlwe}

We next define the intermediate problem $\extlwe$. (This definition is
of an easier problem than the one considered in previous
work~\cite{DBLP:conf/pkc/Alperin-SheriffP12}, which makes our hardness result stronger.)

\begin{definition}\label{def:extlwe}
For $n,m,q,t \ge 1$, $\calZ \subseteq \Z^m$, and a distribution $\chi$ over $\frac{1}{q}\Z^m$, the $\extlwe^t_{n,m,q,\chi,\calZ}$ problem is as follows.
The algorithm gets to choose $\vecz \in \calZ$ and then receives a tuple
\[
(\matA,(\vecb_i)_{i \in [t]},(\inner{\vece_i,\vecz})_{i \in [t]}) \in \T_q^{n \times m} \times (\T_q^m)^t \times ({\textstyle \frac 1 q} \Z)^t.
\]
Its goal is to distinguish between the following two cases.
In the first, $\matA \in \T_q^{n \times m}$ is chosen uniformly,
$\vece_i \in {\textstyle \frac 1 q}\Z^m$ are chosen from $\chi$,
and $\vecb_i = \matA^T \vecs_i + \vece_i \bmod 1$ where $\vecs_i \in \{0,\ldots,q-1\}^n$ are
chosen uniformly.
The second case is identical, except that the $\vecb_i$ are chosen uniformly in $\T_q^m$ independently of everything else. 
\end{definition}


When $t=1$, we omit the superscript $t$. Also, when $\chi$ is $D_{q^{-1}\Z^{m},\alpha}$ for some $\alpha>0$, we replace the
subscript $\chi$ by $\alpha$.
%
We note that a discrete version of $\lwe$ can be defined as a special case of $\extlwe$ by setting $\cZ = \{ 0^m \}$.
%
We next define a measure of quality of sets $\calZ$.



\begin{definition}\label{def:qualityofhintset}
For a real $\xi>0$ and a set $\calZ \subseteq \Z^m$ we say that $\calZ$ is of \emph{quality} $\xi$
if given any $\vecz \in \calZ$, we can efficiently find a unimodular matrix
$\matU \in \Z^{m \times m}$ such that if $\matU'\in \Z^{m \times (m-1)}$ is the matrix obtained
from $\matU$ by removing its leftmost column then all of the columns of $\matU'$ are orthogonal to $\vecz$
and its largest singular value is at most~$\xi$.
\end{definition}

The idea in this definition is that the columns of $\matU'$ 
form a basis of the lattice of integer points that are orthogonal
to $\vecz$, i.e., the lattice $\{\vecb \in \Z^m:
  \inner{\vecb,\vecz} = 0\}$. The quality measures
	how ``short'' we can make this basis.

\begin{claim}\label{clm:qualityofzeroone}
The set $\calZ=\{0,1\}^m$ is of quality $2$.
\end{claim}
\begin{proof}
Let $\vecz \in \calZ$ and assume without loss of generality that its first $k \ge 1$ coordinates are $1$
and the remaining $m-k$ are $0$. Then consider the upper bidiagonal matrix $\matU$ whose diagonal is all $1$s and whose
diagonal above the main diagonal is $(-1,\ldots,-1,0,\ldots,0)$ with $-1$ appearing $k-1$ times. The matrix is clearly
unimodular and all the columns except the first one are orthogonal to $\vecz$. Moreover, by the triangle inequality,
we can bound the operator norm of $\matU$ by the sum of that of the diagonal $1$ matrix and the off-diagonal matrix, both
of which clearly have norm at most $1$.
\end{proof}

\onote{todo: add more quality bounds such as vectors of bounded norm}

\begin{lemma}\label{lem:redtoextlwe}
Let $\calZ \subseteq \Z^m$ be of quality $\xi>0$. Then for any $n,q \ge 1$, $\eps \in (0,1/2)$, and $\alpha,r  \ge (\ln(2m(1+1/\eps))/\pi)^{1/2}/q$,
there is a (transformation) reduction from the first-is-errorless variant of $\lwe_{n,m,q,\alpha}$ to
$\extlwe_{n,m,q,(\alpha^2 \xi^2 + r^2)^{1/2},\calZ}$ that reduces the advantage by at most $33 \eps/2$. 
\end{lemma}
\begin{proof}
We first describe the reduction.
Assume we are asked to provide samples for some $\vecz \in \calZ$.
We compute a unimodular $\matU \in \Z^{m \times m}$ for $\vecz$ as in Definition~\ref{def:qualityofhintset}, and let $\matU' \in \Z^{m \times (m-1)}$ be
the matrix formed by removing the first column of $\matU$.
We then take $m$ samples from the given distribution, resulting in $(\matA, \vecb) \in \T_q^{n \times m} \times (\T_q^{} \times \T^{m-1})$.
We also sample a vector $\vecf$ from the $m$-dimensional continuous Gaussian distribution
$D_{\alpha(\xi^2 \matI - \matU' \matU'^T)^{1/2}}$,
 which is well defined since $\xi^2 \matI - \matU' \matU'^T$ is a positive semidefinite matrix
by our assumption on $\matU$.
The output of the reduction is the tuple
\begin{equation}\label{eq:outputofextlwereduction}
(\matA'= \matA \matU^T,
  \vecb' + \vecc,
	\inner{\vecz , \vecf + \vecc})
	\in \T_q^{n \times m} \times \T_q^m \times {\textstyle{\frac{1}{q}}}\Z,
\end{equation}
where $\vecb'=\matU \vecb + \vecf$, and $\vecc$ is chosen from the discrete
Gaussian distribution $D_{q^{-1}\Z^m - \vecb',r}$ (using Lemma~\ref{lem:gpv}). 

We now prove the correctness of the reduction. Consider first the case that we get valid LWE equations, i.e.,
$\matA$ is uniform in~$\T_q^{n \times m}$ and $\vecb = \matA^T \vecs + \vece \in \T^m$
where $\vecs \in \{0,\ldots,q-1\}^n$ is uniformly chosen, the first coordinate of $\vece \in \R^m$ is~$0$, 
and the remaining $m-1$ coordinates are chosen from~$D_\alpha$.
Since $\matU$ is unimodular, $\matA' = \matA \matU^T$ is uniformly distributed in $\T_q^{n \times m}$ as required.
From now on we condition on an arbitrary $\matA'$ and analyze the distribution of the remaining two
components of~\eqref{eq:outputofextlwereduction}.
Next,
\[
\vecb' =
\matU \vecb + \vecf =
\matA'^T \vecs + \matU \vece + \vecf.
\]
Since $\matU \vece$ is distributed as a continuous Gaussian $D_{\alpha \matU'}$,
the vector~$\matU \vece + \vecf$ is a distributed as a \emph{spherical} continuous 
Gaussian~$D_{\alpha \xi}$.
Moreover, since $\matA'^T \vecs \in \T_q^m$, the coset $q^{-1}\Z^m - \vecb'$ is identical to
$q^{-1}\Z^m - (\matU \vece + \vecf)$, so we can see $\vecc$ as being chosen from
$D_{q^{-1}\Z^m - (\matU \vece + \vecf),r}$. Therefore, by Lemma~\ref{lem:peikertdiscretegauss}
and using that $r \ge \eta_\eps(q^{-1}\Z^m)$ 
by Lemma~\ref{lem:boundonsmoothing},
the distribution of $\matU \vece + \vecf + \vecc$ is within statistical distance $8 \eps$ 
of $D_{q^{-1}\Z^m,(\alpha^2 \xi^2 + r^2)^{1/2}}$. This shows that the second
component in~\eqref{eq:outputofextlwereduction} is also distributed correctly.
Finally, for the third component, by our assumption on $\matU$ and the fact that the
first coordinate of $\vece$ is zero,
\[
\inner{\vecz, \vecf + \vecc} = \inner{\vecz, \matU\vece + \vecf + \vecc},
\]
and so the third component gives the inner product of the noise with $\vecz$, as desired.

We now consider the case where the input is uniform, i.e.,
that $\matA$ is uniform in~$\T_q^{n\times m}$
and $\vecb$ is independent and uniform in $\T_q \times \T^{m-1}$.
We first observe that by Lemma~\ref{lem:smoothingcontinuous}, since $\alpha \ge \eta_{\eps/m}(q^{-1}\Z)$ (by Lemma~\ref{lem:boundonsmoothing}),
the distribution of~$(\matA,\vecb)$ is within statistical distance $\eps/2$
of the distribution of $(\matA, \vece' + \vece)$ where $\vece'$ is
chosen uniformly in~$\T_q^m$, the first coordinate of $\vece$ is zero,
and its remaining $m-1$ coordinates are chosen independently from~$D_\alpha$.
So from now on assume our input is $(\matA,\vece'+\vece)$.
The first component of~\eqref{eq:outputofextlwereduction} is uniform
in $\T_q^{n \times m}$ as before, and moreover, it is clearly independent
of the other two. Moreover, since $\vecb' = \matU \vece' + \matU \vece + \vecf$ and $\matU \vece' \in \T_q^m$,
the coset $q^{-1}\Z^m - \vecb'$ is identical to
$q^{-1}\Z^m - (\matU \vece + \vecf)$, and so $\vecc$ is distributed
identically to the case of a valid LWE equation, and in particular
is independent of $\vece'$. This establishes that the third component
of~\eqref{eq:outputofextlwereduction} is correctly distributed;
moreover, since $\vece'$ is independent of the first and third components,
and $\matU \vece'$ is uniform in $\T_q^m$ (since $\matU$ is unimodular),
we get that the second component is uniform and independent of the other two,
as desired.
\end{proof}

We end this section by stating the standard reduction to the multi-secret ($t \ge 1$) case of extended LWE. 

%The proof is by a straightforward hybrid argument and is omitted.

\onote{removed $\calD$ from the following lemma since we never  use any distribution on secrets but the uniform one}

\begin{lemma}\label{lem:multiinstanceextlwe}
Let $n, m, q, \chi, \cZ$ be as in Definition~\ref{def:extlwe} with $\chi$ efficiently sampleable, and let $t \ge 1$ be an integer. 
Then there is an efficient (transformation) reduction from $\extlwe_{n,m,q,\chi,\cZ}$ to $\extlwe^t_{n,m,q,\chi,\cZ}$ that reduces the advantage by 
a factor of~$t$.
%\dnote{I changed~$\phi$ to $\chi$, for consistency with 
%definition~\ref{def:extlwe}.} 
\end{lemma}

The proof is by a standard hybrid argument. We bring it here for the sake of completeness. We note that the distribution of the secret vector $\vc{s}$ needs to be sampleable but otherwise it plays no role in the proof. The lemma therefore naturally extends to any (sampleable) distribution of $\vc{s}$.

\begin{proof}
Let $\cA$ be an algorithm for $\extlwe^t_{n,m,q,\chi,\cZ}$, let $\vc{z}$ be the vector output by $\cA$ in the first step (note that this is a random variable) and let $H_i$ denote the distribution
\[
\left(\mx{A}, \{\vc{b}_1, \ldots, \vc{b}_i, \vc{u}_{i+1}, \dots, \vc{u}_t\}, \vc{z}, \{\langle\vc{z}, \vc{e}_i\rangle\}_{i\in[t]} \right)~,\]
%\dnote{I added:}
where~$\vc{u}_{i+1}, \dots, \vc{u}_t$ are sampled independently and uniformly in~$\T_q^m$.
Then by definition $\adv[\cA] = \abs{\Pr[\cA(H_0)]- \Pr[\cA(H_t)]}$.

We now describe an algorithm $\cB$ for $\extlwe_{n,m,q,\chi,\cZ}$: First, $\cB$ runs $\cA$ to obtain $\vc{z}$ and sends it to the challenger as its own $\vc{z}$. Then, given an input $(\mx{A}, \vc{d}, \vc{z}, y)$ for $\extlwe_{n,m,q,\chi,\cZ}$, the distinguisher $\cB$ samples $i^* \getsr [t]$, and in addition $\vc{s}_1, \ldots, \vc{s}_{i^*-1}\getsr \bbZ^n_q$, $\vc{e}_1, \ldots, \vc{e}_{i^*-1}, \vc{e}_{i^*+1}, \ldots, \vc{e}_t \getsr \chi^m$\onote{replaced $\phi$ with $\chi$}, $\vc{u}_{i^*+1}, \ldots, \vc{u}_t \getsr \bbT_q^m$. It sets $\vc{b}_i = \mx{A}^T\cdot\vc{s}_i+\vc{e}_i\pmod{1}$, and sends the following to $\cA$:
\[
\left(\mx{A}, \{ \vc{b}_1, \ldots, \vc{b}_{i^*-1}, \vc{d}, \vc{u}_{i^*+1}, \ldots, \vc{u}_t\}, \vc{z}, \{\langle\vc{z},\vc{e}_1\rangle, \ldots,\langle\vc{z}, \vc{e}_{i^*-1}\rangle, y, \langle\vc{z}, \vc{e}_{i^*+1}\rangle, \ldots, \langle\vc{z}, \vc{e}_t\rangle  \} \right)~.
\]
Finally, $\cB$ outputs the same output as $\cA$ did.

Note that when the input to $\cB$ is distributed as $P_0=(\mx{A}, \vc{b}, \vc{z}, \vc{z}^T \cdot \vc{e})$ with~$\vc{b} = \mx{A}^T\cdot\vc{s}+\vc{e}\pmod{1}$, then $\cB$ feeds $\cA$ with exactly the distribution $H_{i^*}$. On the other hand, if the input to $\cB$ is $P_1 = (\mx{A}, \vc{u}, \vc{z}, \vc{z}^T \cdot \vc{e})$ with~$\vc{u} \getsr \bbT_q^m$, then $\cB$ feeds $\cA$ with~$H_{i^*-1}$.

Since $i^*$ is uniform in $[t]$, we get that
\begin{eqnarray*}
t \, \adv[\cB] & = & t \, \abs{\Pr[\cB(P_0)] - \Pr[\cB(P_1)]}\\
  & = &
	\abs[\bigg]{{\sum_{i^*\in[t]} \Pr[\cA(H_{i^*})] - \sum_{i^*\in[t]} \Pr[\cA(H_{i^*-1})]}} \\
	& = &
	\abs{\Pr[\cA(H_{t})]- \Pr[\cA(H_{0})] } \\
	& = & \adv[\cA]~,
\end{eqnarray*}
and the result follows.
\end{proof}




%%% Local Variables:
%%% mode: latex
%%% TeX-master: "lwehard"
%%% End:


\subsection{Reducing to binary secret}\label{sec:theshortsecretreduction}

\begin{lemma}\label{lem:zolwe}
Let $k,n,m,q \in \bbN$, $\epsilon \in (0,1/2)$, and $\delta,\alpha,\beta,\gamma>0$ be such that $n \ge k \log_2 q + 2 \log_2 (1/\delta)$, $\beta \ge \sqrt{2 \ln(2n(1+1/\eps))/\pi}/q$, $\alpha = \sqrt{2n} \beta$, $\gamma = \sqrt{n}\beta$. 
Then there exist three efficient (transformation) reductions to $\zolwe_{n,m,q,\le\alpha}$ from 
$\extlwe^m_{k,n,q,\beta,\{0,1\}^n}$, $\lwe_{k,m,q,\gamma}$, and  $\extlwe^m_{k,n,q,\beta,\{0^n\}}$,
such that if $\cB_1$, $\cB_2$, and $\cB_3$ are the algorithms obtained by applying these reductions (respectively) 
to an algorithm $\cA$, then
\[
\adv[\cA] \le \adv[\cB_1]+\adv[\cB_2]+\adv[\cB_3]+4m\epsilon + \delta \,.
\]
\end{lemma}

Pointing out the trivial (transformation) reduction from $\extlwe^m_{k,n,q,\beta,\{0,1\}^n}$ to $\extlwe^m_{k,n,q,\beta,\{0^n\}}$, the lemma implies the hardness of $\zolwe_{n,m,q,\le\alpha}$ based on the hardness of $\extlwe^m_{k,n,q,\beta,\{0,1\}^n}$ and $\lwe_{k,m,q,\gamma}$.

We note that our proof is actually more general, and holds for any binary distribution of
min-entropy at least $k \log_2 q + 2 \log_2 (1/\delta)$, and not just a uniform binary secret as in the definition of $\zolwe$.


\def\veh{\hat{\vc{e}}}


\begin{proof}
The proof follows by a sequence of hybrids. 
Let $k,n,m,q,\eps, \alpha, \beta, \gamma$ be as in the lemma statement. 
We consider $\vc{z} \getsr \binset^n$ and $\vc{e} \getsr D_{\alpha'}^m$ for $\alpha' = \sqrt{\beta^2\|\vc{z}\|^2+ \gamma^2} \le \sqrt{2n}\beta=\alpha$. In addition, we let $\mx{A}\getsr \bbT_q^{n \times m}$, $\vc{u} \getsr \bbT^m$, and define $\vc{b} \getsd \mx{A}^T\cdot\vc{z}+\vc{e}\pmod{1}$. We consider an algorithm~$\cA$ that distinguishes between $(\mx{A}, \vc{b})$ and $(\mx{A}, \vc{u})$.
%\dnote{I modified the last sentence a bit}

We let $H_0$ denote the distribution $(\mx{A}, \vc{b})$ and $H_1$ the distribution 
\[ H_1 = (\mx{A}, \mx{A}^T\vc{z}-\mx{N}^T\vc{z}+\veh \bmod 1 ),\]
where $\mx{N}\getsr D^{n \times m}_{q^{-1}\bbZ, \beta}$ and $\veh \getsr D_{\gamma}^m$.
Using $\| \vc{z} \| \le \sqrt{n}$ and that $\beta \ge \sqrt{2} \eta_\eps(\Z^n)/q$ (by Lemma~\ref{lem:boundonsmoothing}), it follows by Lemma~\ref{lem:smoothip} that the statistical distance between $-\mx{N}^T\vc{z}+\veh$ and $D^m_{\alpha'}$ is at most $4 m \epsilon$.
It thus follows that
\begin{equation}\label{eq:h0}
\abs{\Pr[\cA(H_0)] - \Pr[\cA(H_1)]} \le 4m\epsilon~.
\end{equation}


We define a distribution $H_2$ as follows. Let $\mx{B} \getsr \bbT_q^{k \times m}$ and $\mx{C} \getsr \bbT_q^{k \times n}$.
Let $\mAh \getsd q\mx{C}^T\cdot \mx{B} + \mx{N}\pmod{1}$. Finally,
\[
H_2 = (\mAh, \mAh^T\cdot\vc{z} - \mx{N}^T\vc{z}+\veh) = (\mAh, q\mx{B}^T\cdot \mx{C}\cdot\vc{z}+\veh)~.
\]

We now argue that there exists an adversary $\cB_1$ for problem $\extlwe^m_{k,n,q,\beta,\{0,1\}^n}$, such that
\begin{equation}\label{eq:h1}
\adv[\cB_1] = \abs{\Pr[\cA(H_1)]-\Pr[\cA(H_2)]}~.
\end{equation}
This is because $H_1, H_2$ can be viewed as applying the same efficient transformation on the distributions $(\mx{C}, \mx{A}, \mx{N}^T\vc{z})$ and $(\mx{C}, \mAh, \mx{N}^T\vc{z})$ respectively. Since distinguishing the latter distributions is exactly the $\extlwe^m_{k,n,q,\beta,\{0,1\}^n}$ problem (where the columns of $q\cdot\mx{B}$ are interpreted as the $m$ secret vectors), the distinguisher $\cB_1$ follows by first applying the aforementioned transformation and then applying~$\cA$.

For the next hybrid, we define $H_3 = (\mAh, \mx{B}^T\cdot \vc{s}+\veh)$, for $\vc{s} \getsr \bbZ_q^k$.
It follows that
\begin{equation}\label{eq:h2}
\abs{\Pr[\cA(H_2)]- \Pr[\cA(H_3)]} \le \delta
\end{equation}
by the leftover hash lemma (see Lemma~\ref{lem:lhl}), since $H_2, H_3$ can be derived from $(\mx{C}, q\mx{C}\cdot\vc{z})$ and $(\mx{C}, \vc{s})$ respectively, whose statistical distance is at most $\delta$.

Our next hybrid makes the second component uniform: $H_4=(\mAh, \vc{u})$. There exists an algorithm $\cB_2$ for $\lwe_{k,m,q,\gamma}$ such that
\begin{equation}\label{eq:h3}
\adv[\cB_2] = \abs{\Pr[\cA(H_3)]- \Pr[\cA(H_4)]}~,
\end{equation}
since $H_3, H_4$ can be computed efficiently from $(\mx{B}, \mx{B}^T\vc{s}+\veh), (\mx{B}, \vc{u})$.

Lastly, we change $\mAh$ back to uniform: $H_5 = (\mx{A}, \vc{u})$. There exists an algorithm $\cB_3$ for $\extlwe^m_{k,n,q,\beta,\{0^n\}}$ such that
\begin{equation}\label{eq:h4}
\adv[\cB_3] = \abs{\Pr[\cA(H_4)]- \Pr[\cA(H_5)]}~.
\end{equation}
Eq.~\eqref{eq:h4} is derived very similarly to Eq.~\eqref{eq:h1}: We notice that $H_4$, $H_5$ can be viewed as applying the same efficient transformation on the distributions $(\mx{C}, \mAh)$ and $(\mx{C}, \mx{A})$ respectively. Since distinguishing the latter distributions is exactly the $\extlwe^m_{k,n,q,\beta,\{0^n\}}$ problem (where the columns of $q\cdot\mx{B}$ are interpreted as the $m$ secret vectors), the distinguisher $\cB_3$ follows by first applying the aforementioned transformation and then applying~$\cA$.

Putting together Eq.~\eqref{eq:h0},~\eqref{eq:h1},~\eqref{eq:h2},~\eqref{eq:h3},~\eqref{eq:h4}, the lemma follows.
\end{proof}












%%% Local Variables:
%%% mode: latex
%%% TeX-master: "lwehard"
%%% End:
